\documentclass[UTF8,a4paper,11pt]{ctexart}
\usepackage[margin=22mm,headheight=14pt,footskip=12mm]{geometry}
\usepackage{amsmath,amssymb,array,longtable,booktabs,tikz,needspace,fancyhdr,hyperref}
\usepackage{microtype}
\hypersetup{colorlinks=false,pdfborder={0 0 0},pdfauthor={Good Learning},bookmarks=false}
\newfontfamily\circledfont{Segoe UI Symbol}
\newcolumntype{P}[1]{>{\raggedright\arraybackslash}p{#1}}
\setlength{\LTleft}{0pt}\setlength{\LTright}{0pt}
\setlength{\parindent}{0pt}\setlength{\parskip}{.65em}
\linespread{1.18}
\raggedbottom
\setcounter{secnumdepth}{0}
\ctexset{section={format=\Large\bfseries\sffamily,beforeskip=1.35em,afterskip=.65em},
subsection={format=\large\bfseries\sffamily,beforeskip=1em,afterskip=.45em}}
\newcommand{\chaptertitle}[1]{\Needspace{8\baselineskip}\noindent{\LARGE\bfseries\sffamily #1}\par\vspace{4mm}\hrule\vspace{6mm}}
\newcommand{\CompactQuestion}[2]{\par\noindent\textbf{#1.}\ #2\par\vspace{.8em}}
\newcommand{\PracticeQuestion}[3]{\par\Needspace{4\baselineskip}\noindent\begin{minipage}{\linewidth}\textbf{#1.}\ #2\par\noindent\rule{0pt}{#3mm}\par\end{minipage}\par}
\newcommand{\CompactQuestionColon}[2]{\par\noindent\textbf{#1：}\ #2\par\vspace{.8em}}
\newcommand{\PracticeQuestionColon}[3]{\par\Needspace{4\baselineskip}\noindent\begin{minipage}{\linewidth}\textbf{#1：}\ #2\par\noindent\rule{0pt}{#3mm}\par\end{minipage}\par}
\pagestyle{fancy}\fancyhf{}\fancyfoot[C]{\thepage}\renewcommand{\headrulewidth}{0pt}
\begin{document}
\thispagestyle{empty}
\noindent{\sffamily\footnotesize 每日课程}\hfill{\sffamily\footnotesize 学习导图}\par
\vspace{7mm}
\noindent{\sffamily\fontsize{25}{30}\selectfont\bfseries 第 8 天 · 第一轮复盘（宽松排版型）}\par
\vspace{4mm}\noindent\rule{\linewidth}{.85pt}\par\vspace{6mm}
\begin{center}
\begin{tikzpicture}[x=1cm,y=1cm,>=stealth,every node/.style={font=\sffamily}]
\path[use as bounding box] (-8.1,-18.7) rectangle (8.1,.4);
\node[font=\sffamily\Large\bfseries] at (0,-.6) {三道闭卷题与错题复做};
\node[font=\sffamily\small,text=black!55] at (0,-1.25) {原定 11 月 22 日 · 11 月 23 日补学};
\node[anchor=north west,font=\sffamily\fontsize{22}{24}\selectfont\bfseries,text=black!50] at (-7.35,-2.35) {01};
\node[anchor=north west,text width=5.8cm,align=left,font=\sffamily\large\bfseries] at (-6.199999999999999,-2.3200000000000003) {最短路};
\node[anchor=north west,text width=5.9cm,align=left,font=\sffamily\small] at (-6.199999999999999,-3.0700000000000003) {先取当前累计费用最小的点\\[.25em]取出终点后倒查前驱};
\draw[black!35,line width=.4pt] (-7.35,-5.050000000000001) -- (-0.2999999999999998,-5.050000000000001);
\node[anchor=north west,font=\sffamily\fontsize{22}{24}\selectfont\bfseries,text=black!50] at (0.85,-2.35) {02};
\node[anchor=north west,text width=5.8cm,align=left,font=\sffamily\large\bfseries] at (2.0,-2.3200000000000003) {梯度下降};
\node[anchor=north west,text width=5.9cm,align=left,font=\sffamily\small] at (2.0,-3.0700000000000003) {每轮在当前位置算导数\\[.25em]用学习率缩放后更新参数};
\draw[black!35,line width=.4pt] (0.85,-5.050000000000001) -- (7.8999999999999995,-5.050000000000001);
\node[anchor=north west,font=\sffamily\fontsize{22}{24}\selectfont\bfseries,text=black!50] at (0.85,-7.7) {03};
\node[anchor=north west,text width=5.8cm,align=left,font=\sffamily\large\bfseries] at (2.0,-7.67) {回归树};
\node[anchor=north west,text width=5.9cm,align=left,font=\sffamily\small] at (2.0,-8.42) {各叶用组内均值预测\\[.25em]比较总 SSE 或全体 MSE};
\draw[black!35,line width=.4pt] (0.85,-10.4) -- (7.8999999999999995,-10.4);
\node[anchor=north west,font=\sffamily\fontsize{22}{24}\selectfont\bfseries,text=black!50] at (-7.35,-7.7) {04};
\node[anchor=north west,text width=5.8cm,align=left,font=\sffamily\large\bfseries] at (-6.199999999999999,-7.67) {平方误差};
\node[anchor=north west,text width=5.9cm,align=left,font=\sffamily\small] at (-6.199999999999999,-8.42) {总 SSE 除以总人数\\[.25em]按人数加权各叶 MSE 可得同值};
\draw[black!35,line width=.4pt] (-7.35,-10.4) -- (-0.2999999999999998,-10.4);
\node[anchor=north west,font=\sffamily\fontsize{22}{24}\selectfont\bfseries,text=black!50] at (-7.35,-13.05) {05};
\node[anchor=north west,text width=5.8cm,align=left,font=\sffamily\large\bfseries] at (-6.199999999999999,-13.020000000000001) {梯度提升树};
\node[anchor=north west,text width=5.9cm,align=left,font=\sffamily\small] at (-6.199999999999999,-13.770000000000001) {残差均值是新树的叶值\\[.25em]乘学习率后加回旧预测};
\draw[black!35,line width=.4pt] (-7.35,-15.75) -- (-0.2999999999999998,-15.75);
\node[anchor=north west,font=\sffamily\fontsize{22}{24}\selectfont\bfseries,text=black!50] at (0.85,-13.05) {06};
\node[anchor=north west,text width=5.8cm,align=left,font=\sffamily\large\bfseries] at (2.0,-13.020000000000001) {验证与测试};
\node[anchor=north west,text width=5.9cm,align=left,font=\sffamily\small] at (2.0,-13.770000000000001) {验证集比较候选\\[.25em]锁定方案后做最终测试};
\draw[black!35,line width=.4pt] (0.85,-15.75) -- (7.8999999999999995,-15.75);
\draw[->,black!55,line width=.8pt] (-.12,-3.7) -- (.55,-3.7);
\draw[->,black!55,line width=.8pt] (7.85,-5.45) -- (7.85,-7.2);
\draw[->,black!55,line width=.8pt] (.52,-9.05) -- (-.15,-9.05);
\draw[->,black!55,line width=.8pt] (-7.85,-10.8) -- (-7.85,-12.55);
\draw[->,black!55,line width=.8pt] (-.12,-14.4) -- (.55,-14.4);
\node[font=\sffamily\small,text=black!55] at (0,-18.1) {25 分钟闭卷 · 20 分钟订正 · 15 分钟补漏或休息};
\end{tikzpicture}
\end{center}
\clearpage

\chaptertitle{第 8 天 · 第一轮复盘}
\textbf{原计划 2026 年 11 月 22 日｜补学 2026 年 11 月 23 日｜K03–K10｜60 分钟}

前 25 分钟闭卷做 P1、P2、P3，分别用 8、8、9 分钟。每题写出中间步骤，卡住就先标记。接着用 20 分钟看相应例题、核对答案，再做 R1、R2、R3。最后 15 分钟按卡点挑一至两道填步题或变式题复做，也可以休息。

\textbf{验收。} P1 写出取点顺序、路径和总代价。P2 每轮用当前值算导数并更新。P3 写出分组、叶预测、分裂误差和选择理由。三类至少两类独立写对主步骤。把题号、算式、卡住的一步和实际用时发回来，便于判断哪些需要复做。

\textbf{题目来源。} P3 按备考计划和知识清单所记的“计算 18 连续价格分裂”考点新编。R1 至 R3 来自本仓库错题 15、18、19 的复做题。其余题为新编。

\section{独立练习 · 先用 25 分钟闭卷做}
\PracticeQuestion{P1}{从 S 到 G 的有向路及路费为 S→A 2、S→B 5、A→B 1、A→G 7、B→G 2。用 Dijkstra 写出取点顺序、最便宜的路线及总路费。指出第一次发现 G 时和最终确定 G 时各花多少钱。}{35}
\PracticeQuestion{P2}{损失 $J(w)=(w-2)^2+1$，初值 $w_0=6$，学习率 $\eta=1/4$。做两步梯度下降，列出每步的导数和更新后的 $w$。算 $J(w_0)$、$J(w_2)$，检查损失是否下降。}{35}
\PracticeQuestion{P3}{四套房的特征 $x$ 依次是 $1,2,3,4$，价格依次是 $2,4,8,12$。预测连续价格，候选阈值只有 $1.5$ 和 $2.5$，约定 $x\le t$ 走左边。按平方误差准则比较两个阈值。列出每个候选的左右组、叶预测、总 SSE 和全体 MSE，再选阈值。若题目把“分裂前 MSE 减分裂后 MSE”称为增益，求胜出阈值的增益。}{35}
\clearpage
\section{单元 1 · 路径题}
Dijkstra 处理边权非负的图。给起点记 0，给还没有发现的地点记无穷大。每轮从尚未确定的地点中，取当前累计花费最小的一个。经它走到邻居，只有总花费变小时才更新邻居的费用和前驱。从队列取出终点时，可以沿前驱倒着找路。

\subsection{例题 · 什么时候能确定终点}
有向路为 S→A 2、S→B 4、A→G 5、B→G 1。

\begin{enumerate}
\item 取 S，得到 A 费 2、B 费 4。先取 A，因为 2 小于 4。
\item 经 A 到 G 花 $2+5=7$。此时 B 的 4 仍比 G 的 7 小，先取 B。
\item 经 B 到 G 花 $4+1=5$，比 7 便宜，更新 G 的费用和前驱。
\item 现在取 G。倒查得到 S→B→G，顺着加一遍是 $4+1=5$。
\end{enumerate}
\subsection{举一反三}
\PracticeQuestion{A1}{S→A 1、S→B 4、A→G 6、B→G 2。经 A 已记下 G 费 7。取 B 后，G 的费用和前驱应填什么？}{23}
\PracticeQuestion{A2}{沿用 A1，把 B→G 的路费改为 8。写取点顺序和最终路径，说明哪次更新会被跳过。}{26}
\clearpage
\section{单元 2 · 梯度下降}
一元函数在当前点的导数告诉你损失沿 $w$ 增大方向的变化快慢。先把当前的 $w$ 代入导数，再用 $w_{k+1}=w_k-\eta J'(w_k)$ 更新。导数为负时，更新会增大 $w$。学习率过大可能跨过最低点，算完可以代回原损失检查。

\subsection{例题 · 每步用新的位置}
设 $J(w)=(w-3)^2+2$，$w_0=1$，$\eta=1/4$。

\begin{enumerate}
\item 先求 $J'(w)=2(w-3)$。常数 2 的导数为 0。
\item 在 $w_0=1$ 处，导数为 $-4$，所以 $w_1=1-(1/4)(-4)=2$。
\item 在 $w_1=2$ 处重算，导数为 $-2$，所以 $w_2=2-(1/4)(-2)=2.5$。
\item 代回整个损失，$J(1)=6$，$J(2.5)=2.25$。这两步让损失下降。
\end{enumerate}
\subsection{举一反三}
\PracticeQuestion{B1}{沿用例题函数和学习率，若某轮当前值为 4，此轮导数和更新值是什么？}{24}
\PracticeQuestion{B2}{保持例题初值，把学习率改为 1。做两步并比较前后损失，指出改变学习率带来的结果。}{34}
\clearpage
\section{单元 3 · 回归树分裂}
在平方误差准则下，每个叶子用组内价格的均值做预测。SSE 是每个样本的误差平方后相加，MSE 是 SSE 除以样本数。比较同一批样本的候选分裂时，可以直接比较总 SSE，因为 MSE 的分母相同。若先各自算出左右叶 MSE，要按各叶样本所占比例加权。

\subsection{例题 · 两个叶子人数不同时}
$x=1,2,3$ 的价格依次为 $2,4,8$，比较阈值 $1.5$ 和 $2.5$。约定 $x\le t$ 走左边。

\begin{enumerate}
\item 阈值 $1.5$ 分出 $\{2\}$、$\{4,8\}$。均值为 2、6。两叶 SSE 分别为 0 和 $(4-6)^2+(8-6)^2=8$，合计 8。
\item 阈值 $2.5$ 分出 $\{2,4\}$、$\{8\}$。均值为 3、8。两叶 SSE 分别为 2 和 0，合计 2，因此选 $2.5$。
\item 胜出分裂的 MSE 为 $2/3$。也可先算两叶 MSE 为 1、0，再算 $(2/3)\times1+(1/3)\times0=2/3$。两种算法数的是同一批误差。
\end{enumerate}
\subsection{举一反三}
\PracticeQuestion{C1}{某次分裂左叶 1 人、SSE 为 0，右叶 4 人、SSE 为 12。右叶 MSE 是 3，补出分裂后的全体 MSE。}{27}
\PracticeQuestion{C2}{把例题价格改为 $2,7,8$。重算两个候选的叶预测和总 SSE。最佳阈值有什么变化？}{35}
\clearpage
\section{单元 4 · 订正后复做}
选模型时，用训练集拟合各候选，用验证集比较。锁定方案后用测试集做最终评估。

平方损失下，梯度提升树本轮先算每个样本的残差 $r=y-F_{\text{旧}}$。新树在每个叶上输出该叶的残差均值 $h$。把 $\eta h$ 加到旧预测上，才得到模型新预测。

\subsection{例题 · 区分叶值、修正量和预测}
真实值为 $2,6$，旧预测都是 3，学习率 $\eta=1/2$。假设本轮两个样本在同一个叶中。

\begin{enumerate}
\item 用真实值减旧预测，得残差 $(-1,3)$。
\item 同叶取均值，得叶值 $h=1$。两条样本共享这一个叶值。
\item 乘学习率，修正量为 $0.5$。加回旧预测，两条新预测都是 $3.5$。
\item 下一轮重新算残差，得 $(-1.5,2.5)$。
\end{enumerate}
\subsection{举一反三}
\PracticeQuestion{R1}{对应错题 15。$\lambda$ 取 $0,0.1,1$ 时，验证错误率依次为 $16\%,9\%,12\%$。选出 $\lambda$，写出什么时候使用测试集。}{18}
\PracticeQuestion{R2}{对应错题 18。分裂后左叶有 2 人、SSE 为 2，右叶有 3 人、SSE 为 12。分别写两叶 MSE，再用两种方法算全体 MSE。}{26}
\PracticeQuestion{R3}{对应错题 19。真实值为 $2,4,10$，初始预测都是 4，学习率为 $1/2$。每轮固定前两条样本同叶，最后一条单独成叶。写第一轮的残差、左右叶值、修正量和新预测。若有余力，用新预测再做第二轮。}{42}
\clearpage
\section{核对答案}
\textbf{P1.} 取 S，暂记 A 为 2、B 为 5。取 A，把 B 更新为 $2+1=3$，首次记 G 为 $2+7=9$。取 B，把 G 更新为 $3+2=5$，最后取 G。顺序 S、A、B、G，路径 S→A→B→G，总代价 5。

\textbf{P2.} $J'(w)=2(w-2)$。第一步导数为 8，$w_1=6-8/4=4$。第二步在 4 求导，得 4，$w_2=4-4/4=3$。完整损失从 $J(6)=17$ 降到 $J(3)=2$。

\textbf{P3.} 阈值 $1.5$ 分出价格 $\{2\}$、$\{4,8,12\}$。叶预测是 2、8，总 SSE 为 $0+16+0+16=32$，MSE 为 8。阈值 $2.5$ 分出 $\{2,4\}$、$\{8,12\}$。叶预测是 3、10，两边分别贡献 2 和 8，总 SSE 为 10，MSE 为 2.5，因此选 $2.5$。分裂前均值为 6.5，MSE 为 $(20.25+6.25+2.25+30.25)/4=14.75$，胜出分裂的 MSE 降低量为 $14.75-2.5=12.25$。

\textbf{A1.} 取 B 后，经 B 到 G 花 $4+2=6$，比 7 小，更新费用为 6、前驱为 B。

\textbf{A2.} 顺序 S、A、B、G。经 B 到 G 的费用变成 12，高于已记下的 7。跳过这次更新，保留路径 S→A→G，总代价 7。

\textbf{B1.} $J'(4)=2$，更新为 $4-(1/4)\times2=3.5$。

\textbf{B2.} 第一步导数 $-4$，$w_1=5$。第二步导数 4，$w_2=1$。三次位置的损失都是 6。这个步长让 $w$ 在最低点 3 的两边来回跳。

\textbf{C1.} 全体 MSE 为 $(1/5)\times0+(4/5)\times3=12/5$，也等于总 SSE 12 除以总人数 5。

\textbf{C2.} 阈值 $1.5$ 的叶预测是 2、7.5，总 SSE 为 $0+0.5=0.5$。阈值 $2.5$ 的叶预测是 4.5、8，总 SSE 为 $12.5+0=12.5$。改选 $1.5$，因为它分出的两组里，价格分别更接近各自的均值。

\textbf{R1.} 选 $\lambda=0.1$，因为验证错误率 9\% 最低。锁定模型及处理流程后，才用测试集做最终评估。

\textbf{R2.} 左、右 MSE 分别为 $2/2=1$ 和 $12/3=4$。按人数加权得 $(2/5)\times1+(3/5)\times4=14/5$。也可先加总两叶 SSE，再除以人数，得 $(2+12)/5=14/5$。

\textbf{R3.} 第一轮残差 $(-2,0,6)$，左右叶值 $-1,6$，乘学习率后修正量 $-0.5,3$，加回旧预测得到 $(3.5,3.5,7)$。第二轮用第一轮的新预测算残差 $(-1.5,0.5,3)$，左右叶值 $-0.5,3$，修正量 $-0.25,1.5$，得到预测 $(3.25,3.25,8.5)$。

\end{document}
